\documentclass[11pt,a4paper,oneside]{article}

\usepackage[T1]{fontenc}
\usepackage[utf8]{inputenc}
\usepackage[polish]{babel}
\usepackage{lmodern}
\usepackage{microtype}
\usepackage[a4paper,margin=2.25cm]{geometry}
\usepackage{amsmath,amssymb,mathtools}
\usepackage{bm}
\usepackage{booktabs,longtable,array,tabularx}
\usepackage{enumitem}
\usepackage{xcolor}
\usepackage{hyperref}
\usepackage[nameinlink,noabbrev]{cleveref}
\usepackage{graphicx}
\usepackage{tikz}
\usetikzlibrary{arrows.meta,positioning,fit,shapes.geometric}
\usepackage{listings}
\usepackage{fancyhdr}
\usepackage{lastpage}
\usepackage{siunitx}

\definecolor{codegray}{RGB}{245,245,245}
\definecolor{darkblue}{RGB}{20,55,100}
\definecolor{darkgreen}{RGB}{25,105,75}
\definecolor{darkred}{RGB}{145,40,40}

\hypersetup{
  colorlinks=true,
  linkcolor=darkblue,
  urlcolor=darkblue,
  citecolor=darkgreen,
  pdftitle={Godzinowy wielozadaniowy model prognozowania jakości powietrza i pogody},
  pdfauthor={Edward Szczypka}
}

\lstdefinestyle{code}{
  backgroundcolor=\color{codegray},
  basicstyle=\ttfamily\small,
  breaklines=true,
  frame=single,
  columns=fullflexible,
  showstringspaces=false,
  keywordstyle=\color{darkblue},
  commentstyle=\color{darkgreen},
  stringstyle=\color{darkred}
}

\setlength{\headheight}{14pt}
\pagestyle{fancy}
\fancyhf{}
\lhead{Smog AI --- model godzinowy}
\rhead{Edward Szczypka}
\cfoot{\thepage/\pageref{LastPage}}

\newcommand{\PMten}{\mathrm{PM}_{10}}
\newcommand{\PMtwo}{\mathrm{PM}_{2.5}}
\newcommand{\Rplus}{\mathbb{R}_{+}}
\newcommand{\E}{\mathbb{E}}
\newcommand{\Prob}{\mathbb{P}}
\newcommand{\ind}{\mathbf{1}}
\newcommand{\Loss}{\mathcal{L}}
\newcommand{\X}{\mathbf{X}}
\newcommand{\Y}{\mathbf{Y}}
\newcommand{\thetaT}{\theta_T}
\newcommand{\thetaR}{\theta_R}
\newcommand{\thetaPM}{\theta_{\mathrm{PM}}}

\title{\textbf{Godzinowy, warunkowany horyzontem model prognozowania}\\
\textbf{PM10, PM2.5, temperatury i opadu}\\[0.4em]
\large Specyfikacja matematyczna, walidacyjna i wdrożeniowa}
\author{Edward Szczypka\\\href{mailto:edward@szczypka.guru}{edward@szczypka.guru}}
\date{\today}

\begin{document}
\maketitle

\begin{abstract}
Dokument opisuje docelowy model systemu Smog AI, w którym prognoza jest
wyznaczana dla \emph{dokładnego czasu docelowego}, a nie przez wybór najbliższego
z trzech dyskretnych horyzontów 6, 12 i 24 godzin. System prognozuje cztery
wielkości:
\[
  \PMten,\qquad \PMtwo,\qquad T,\qquad R,
\]
gdzie $T$ oznacza temperaturę, a $R$ sumę opadu w jawnym okresie akumulacji
$\Delta_R$ kończącym się w czasie docelowym. Domyślnie $\Delta_R=6$ godzin,
zgodnie z bieżącym kontraktem danych IMGW; system nie tworzy sztucznej
dezagregacji na wartości godzinowe. Obliczenia, trening, predykcja i
interpolacja przestrzenna są
wykonywane lokalnie. DigitalOcean Spaces przechowuje wersjonowane artefakty,
a DigitalOcean App Platform wyłącznie odczytuje i prezentuje gotowe wyniki.
\end{abstract}

\tableofcontents
\newpage

\section{Cel zmiany}

Dotychczasowe modele były trenowane oddzielnie dla horyzontów
\[
  h\in\{6,12,24\}.
\]
Taki model jest poprawny wyłącznie dla czasu
\[
  u=t+h,
\]
gdzie $t$ jest czasem bazowym prognozy. Nie jest poprawne traktowanie modelu
$h=6$ jako prognozy dla $h=22$ ani ciche wybieranie powierzchni o najbliższym
czasie docelowym.

W wersji godzinowej horyzont staje się zmienną objaśniającą:
\[
  h\in\{1,2,\ldots,H\},\qquad H=48.
\]
Model uczy się zależności pomiędzy stanem układu w chwili $t$, dokładnym
horyzontem $h$, cechami czasu docelowego $t+h$ oraz wartością prognozowaną.

\section{Notacja}

Niech:
\begin{itemize}[nosep]
  \item $\mathcal{S}_A$ oznacza zbiór stacji jakości powietrza GIOŚ;
  \item $\mathcal{S}_W$ oznacza zbiór stacji meteorologicznych IMGW;
  \item $m(s)\in\mathcal{S}_W$ oznacza stację IMGW dopasowaną do stacji
        $s\in\mathcal{S}_A$;
  \item $t$ oznacza czas bazowy danych dostępnych przy generowaniu prognozy;
  \item $h$ oznacza horyzont w godzinach;
  \item $u=t+h$ oznacza dokładny czas docelowy;
  \item $P^{10}_{s,t}$ oraz $P^{2.5}_{s,t}$ oznaczają stężenia
        $\PMten$ i $\PMtwo$;
  \item $T_{m,t}$ oznacza temperaturę;
  \item $R^{(\Delta_R)}_{m,t}$ oznacza sumę opadu w okresie
        $[t-\Delta_R,t]$; domyślnie $\Delta_R=6$ godzin.
\end{itemize}

Wektor docelowy dla stacji $s$, czasu bazowego $t$ i horyzontu $h$ ma postać:
\begin{equation}
  \Y_{s,t,h} =
  \begin{bmatrix}
    P^{10}_{s,t+h}\\
    P^{2.5}_{s,t+h}\\
    T_{m(s),t+h}\\
    R^{(\Delta_R)}_{m(s),t+h}
  \end{bmatrix}.
  \label{eq:target-vector}
\end{equation}

\section{Dlaczego nie stosujemy globalnego wielomianu przez punkty 6/12/24 h}

Wielomian drugiego stopnia przechodzący przez trzy punkty
\[
  (6,\widehat y_6),\quad (12,\widehat y_{12}),\quad (24,\widehat y_{24})
\]
jest jednoznacznie wyznaczony, lecz jego gładkość jest wyłącznie własnością
algebraiczną. Nie uwzględnia on zmiany temperatury, opadu, wiatru, pory dnia
ani emisji. Może więc interpolować gładko, ale fizycznie niewiarygodnie.

Globalny wielomian nie jest modelem prognozowania. W projekcie dopuszczamy
interpolację wielomianową wyłącznie na końcowym, lokalnym etapie:
\begin{itemize}[nosep]
  \item pomiędzy dwiema lub kilkoma \emph{już wyznaczonymi prognozami
        godzinowymi};
  \item bez ekstrapolacji poza zakres dostępnych prognoz;
  \item z jawnym oznaczeniem metody.
\end{itemize}

Domyślną metodą dla czasu niecałkowitego jest PCHIP, czyli przedziałowy
wielomian Hermite'a zachowujący lokalny kształt. Dla zapytania o pełną godzinę
interpolacja czasowa nie jest potrzebna.

\section{Architektura modelu}

\subsection{Model kaskadowy}

Ze względu na różny charakter zmiennych nie stosujemy jednej wspólnej funkcji
straty dla wszystkich czterech wyjść. Domyślną architekturą jest kaskada:

\begin{enumerate}
  \item model temperatury;
  \item dwuetapowy model opadu;
  \item modele $\PMten$ i $\PMtwo$, które korzystają z prognozowanej
        temperatury i prognozowanego opadu dla czasu docelowego.
\end{enumerate}

\begin{figure}[ht]
\centering
\resizebox{\textwidth}{!}{%
\begin{tikzpicture}[
  node distance=10mm and 14mm,
  block/.style={draw, rounded corners, align=center, minimum width=37mm,
                minimum height=10mm, fill=blue!5},
  data/.style={draw, rounded corners, align=center, minimum width=37mm,
               minimum height=10mm, fill=green!7},
  arrow/.style={-{Latex[length=2.2mm]}, thick}
]
\node[data] (x) {Cechy dostępne\\w czasie $t$};
\node[block, right=of x, yshift=13mm] (temp) {Model temperatury\\$\widehat T_{t+h}$};
\node[block, right=of x, yshift=-13mm] (rain) {Model opadu\\$\widehat\pi_{t+h},\widehat A_{t+h}$};
\node[block, right=30mm of temp, yshift=-13mm] (pm) {Modele jakości powietrza\\
$\widehat{\PMten}_{t+h}$, $\widehat{\PMtwo}_{t+h}$};
\node[data, right=of pm] (out) {Prognoza godzinowa\\i przedziały niepewności};

\draw[arrow] (x) -- node[above,sloped] {$h,\phi(t+h)$} (temp);
\draw[arrow] (x) -- node[below,sloped] {$h,\phi(t+h)$} (rain);
\draw[arrow] (x) -- (pm);
\draw[arrow] (temp) -- (pm);
\draw[arrow] (rain) -- (pm);
\draw[arrow] (pm) -- (out);
\draw[arrow] (temp.east) |- (out.north);
\draw[arrow] (rain.east) |- (out.south);
\end{tikzpicture}%
}
\caption{Kaskadowa architektura modelu godzinowego.}
\end{figure}

\subsection{Cechy czasu docelowego}

Czas docelowy $u=t+h$ jest kodowany cyklicznie:
\begin{align}
  \phi_{\mathrm{hour},1}(u) &= \sin\left(2\pi\frac{\operatorname{hour}(u)}{24}\right),\\
  \phi_{\mathrm{hour},2}(u) &= \cos\left(2\pi\frac{\operatorname{hour}(u)}{24}\right),\\
  \phi_{\mathrm{dow},1}(u) &= \sin\left(2\pi\frac{\operatorname{dow}(u)}{7}\right),\\
  \phi_{\mathrm{dow},2}(u) &= \cos\left(2\pi\frac{\operatorname{dow}(u)}{7}\right),\\
  \phi_{\mathrm{year},1}(u) &= \sin\left(2\pi\frac{\operatorname{doy}(u)}{365.2425}\right),\\
  \phi_{\mathrm{year},2}(u) &= \cos\left(2\pi\frac{\operatorname{doy}(u)}{365.2425}\right).
\end{align}

Horyzont $h$ jest wprowadzany jawnie oraz opcjonalnie przez bazę funkcji:
\[
  1,\quad h,\quad h^2,\quad \log(1+h),\quad \sqrt{h}.
\]
Nie oznacza to globalnej regresji wielomianowej po wyniku. Są to jedynie cechy,
z których nieliniowy model może skorzystać, jeżeli poprawiają walidację.

\section{Model temperatury}

Model temperatury ma postać:
\begin{equation}
  \widehat T_{s,t,h}
  =
  f_T\!\left(
    \X^{\mathrm{weather}}_{m(s),t},
    h,
    \phi(t+h),
    z_{m(s)},
    \operatorname{lat}_{m(s)},
    \operatorname{lon}_{m(s)}
  \right).
  \label{eq:temperature-model}
\end{equation}

Wektor $\X^{\mathrm{weather}}$ zawiera między innymi:
\[
  T_{t-1},T_{t-3},T_{t-6},T_{t-12},T_{t-24},
\]
średnie kroczące, trend temperatury, wilgotność, ciśnienie, wiatr i opad.

Domyślnym estymatorem jest histogramowy gradient boosting. Modele bazowe:
\begin{align}
  \widehat T^{\mathrm{pers}}_{t+h} &= T_t,\\
  \widehat T^{\mathrm{season}}_{t+h} &=
  \operatorname{median}\{T_{\tau}:
  \operatorname{hour}(\tau)=\operatorname{hour}(t+h),
  \operatorname{month}(\tau)=\operatorname{month}(t+h)\}.
\end{align}

Dla estymacji punktowej stosujemy MAE lub stratę kwadratową, a dla przedziałów
predykcji osobne modele kwantylowe:
\begin{equation}
  \rho_\tau(e)=
  \begin{cases}
    \tau e, & e\ge 0,\\
    (\tau-1)e, & e<0.
  \end{cases}
\end{equation}

\section{Dwuetapowy model opadu}

Opad ma rozkład z dużą masą w zerze i silną prawostronną asymetrią. Dlatego
stosujemy model typu \emph{hurdle}.

Niech $\delta\ge 0$ będzie progiem mierzalnego opadu, np.
$\delta=\SI{0.1}{\milli\meter}$. Definiujemy:
\begin{equation}
  Z_{s,t,h}=\ind\{R^{(\Delta_R)}_{m(s),t+h}>\delta\}.
\end{equation}

\subsection{Prawdopodobieństwo wystąpienia opadu}

\begin{equation}
  \widehat\pi_{s,t,h}
  =
  \Prob(Z_{s,t,h}=1\mid \X_{s,t},h,\phi(t+h))
  =
  \sigma\!\left(f_\pi(\X_{s,t},h,\phi(t+h))\right),
\end{equation}
gdzie $\sigma$ oznacza funkcję logistyczną. Strata:
\begin{equation}
  \Loss_{\mathrm{occ}}
  =
  -\sum_i
  \left[
    z_i\log\widehat\pi_i+
    (1-z_i)\log(1-\widehat\pi_i)
  \right].
\end{equation}

\subsection{Warunkowa wielkość opadu}

Dla obserwacji dodatnich:
\begin{equation}
  A_{s,t,h}=R^{(\Delta_R)}_{m(s),t+h}\mid Z_{s,t,h}=1.
\end{equation}

Modelujemy $\E[A\mid Z=1,\X,h]$ regresją dodatnią, np. z funkcją łączącą
logarytmiczną lub przez transformację $\log(1+A)$. Prognoza średniego opadu:
\begin{equation}
  \widehat R_{s,t,h}
  =
  \widehat\pi_{s,t,h}\,\widehat\mu_{A,s,t,h}.
  \label{eq:rain-expectation}
\end{equation}

Oprócz wartości oczekiwanej aplikacja pokazuje:
\[
  \widehat\pi_{s,t,h},
  \qquad
  \widehat\mu_{A,s,t,h},
  \qquad
  \widehat R_{s,t,h}.
\]

\paragraph{Semantyka okresu akumulacji.}
Wartość $\widehat R_{s,t,h}$ nie oznacza natężenia w mm/h, lecz akumulację
w przedziale długości $\Delta_R$ kończącym się w $t+h$. Okres jest częścią
kontraktu modelu, artefaktu i odpowiedzi API. Zmiana źródła danych może ustawić
$\Delta_R=1$ godzinę, ale nie wolno bezpośrednio dzielić pomiaru sześciogodzinnego
przez sześć.
Dzięki temu wartość \SI{0.4}{\milli\meter} nie jest mylnie interpretowana jako
pewny, ciągły opad.

\section{Modele PM10 i PM2.5}

Dla $k\in\{10,2.5\}$:
\begin{equation}
  \widehat P^{k}_{s,t,h}
  =
  f_k\!\left(
    \X^{\mathrm{air}}_{s,t},
    \X^{\mathrm{weather}}_{m(s),t},
    h,
    \phi(t+h),
    \widehat T_{s,t,h},
    \widehat R_{s,t,h},
    \widehat\pi_{s,t,h},
    \operatorname{lat}_s,
    \operatorname{lon}_s
  \right).
  \label{eq:pm-model}
\end{equation}

Cechy jakości powietrza obejmują opóźnienia, średnie kroczące, tempo zmian
oraz informacje z pobliskich stacji. Prognozowana temperatura i opad są
cechami czasu docelowego, a nie obserwacjami pobranymi z przyszłości.

\section{Ochrona przed wyciekiem informacji}

Najważniejszym niezmiennikiem jest:
\[
  \text{każda cecha użyta w prognozie utworzonej w chwili }t
  \text{ musi być dostępna w chwili }t.
\]

Nie wolno trenować modelu PM na rzeczywistej temperaturze z chwili $t+h$,
jeżeli w produkcji dostępna będzie tylko jej prognoza. Zastosowana zostanie
chronologiczna predykcja krzyżowa:

\begin{enumerate}
  \item dzielimy czasy bazowe na uporządkowane bloki
        $I_1<I_2<\cdots<I_K$;
  \item dla bloku $I_j$ model pogodowy jest trenowany wyłącznie na blokach
        wcześniejszych;
  \item generujemy prognozy
        $\widetilde T_i,\widetilde R_i,\widetilde\pi_i$ dla $i\in I_j$;
  \item modele PM są trenowane na tych prognozach, nie na rzeczywistej
        pogodzie z przyszłości.
\end{enumerate}

Formalnie:
\begin{equation}
  \widetilde T_i =
  f_T^{(-j)}(\X_i,h_i),\qquad i\in I_j,
\end{equation}
gdzie $f_T^{(-j)}$ nie widziało żadnej obserwacji z bloku $I_j$ ani z bloków
późniejszych.

\section{Podział chronologiczny i ocena}

Nie stosujemy losowego podziału. Używamy walidacji typu rolling origin lub
expanding window:
\[
  \underbrace{[1,\ldots,t_1]}_{\text{trening}}
  \rightarrow
  \underbrace{[t_1+1,\ldots,t_2]}_{\text{walidacja}},
\]
następnie powiększamy zbiór treningowy.

Metryki są obliczane:
\begin{itemize}[nosep]
  \item osobno dla każdej zmiennej;
  \item osobno dla każdego horyzontu $h=1,\ldots,48$;
  \item według stacji, pory dnia, sezonu i warunków pogodowych;
  \item łącznie i względem modelu bazowego.
\end{itemize}

Dla temperatury i PM:
\begin{align}
  \operatorname{MAE} &=
  \frac{1}{n}\sum_{i=1}^{n}|y_i-\widehat y_i|,\\
  \operatorname{RMSE} &=
  \sqrt{\frac{1}{n}\sum_{i=1}^{n}(y_i-\widehat y_i)^2},\\
  R^2 &=1-\frac{\sum_i(y_i-\widehat y_i)^2}
                   {\sum_i(y_i-\overline y)^2}.
\end{align}

Dla opadu:
\begin{itemize}[nosep]
  \item Brier score i ROC-AUC dla wystąpienia opadu;
  \item MAE/RMSE dla dodatniej wielkości opadu;
  \item MAE i Tweedie deviance dla całkowitej prognozy opadu;
  \item trafność progów, np. $R>\SI{1}{\milli\meter}$.
\end{itemize}

\section{Interpolacja czasowa}

Model generuje prognozy dla każdej pełnej godziny:
\[
  h=1,2,\ldots,48.
\]
Dla zapytania o pełną godzinę wybór jest dokładny:
\[
  \widehat y(u)=\widehat y_{h},\qquad u=t+h.
\]

Dla zapytania np. o 17:30 dopuszczamy interpolację pomiędzy prognozami
17:00 i 18:00. Domyślnie:
\begin{itemize}[nosep]
  \item interpolacja liniowa przy dwóch sąsiednich punktach;
  \item PCHIP przy wystarczającej liczbie sąsiednich prognoz;
  \item brak ekstrapolacji poza $[t+1,t+48]$.
\end{itemize}

API zwraca:
\begin{lstlisting}[style=code]
{
  "requested_target_time": "...",
  "selected_target_time": "...",
  "forecast_origin_time": "...",
  "horizon_hours": 22.0,
  "exact_time_match": true,
  "temporal_method": "direct_hourly_model"
}
\end{lstlisting}

\section{Interpolacja przestrzenna}

Dla ustalonego czasu docelowego $u$ i zmiennej $y$ otrzymujemy prognozy
punktowe $\widehat y_i(u)$ dla stacji $i=1,\ldots,n$. W wariancie IDW:
\begin{align}
  w_i(x) &= \frac{1}{(d(x,x_i)+\varepsilon)^p},\\
  \widehat y(x,u)
  &=
  \frac{\sum_{i\in\mathcal N_k(x)}w_i(x)\widehat y_i(u)}
       {\sum_{i\in\mathcal N_k(x)}w_i(x)}.
  \label{eq:idw}
\end{align}

Dla opadu interpolujemy oddzielnie:
\[
  \widehat\pi(x,u)
  \quad\text{oraz}\quad
  \widehat\mu_A(x,u),
\]
a następnie:
\[
  \widehat R(x,u)=\widehat\pi(x,u)\widehat\mu_A(x,u).
\]
Ogranicza to sztuczne rozsmarowywanie niewielkiego opadu na cały kraj.

Dla temperatury możliwa jest opcjonalna korekta wysokościowa:
\begin{equation}
  \widehat T^{\,*}_i(u;x)
  =
  \widehat T_i(u)+\gamma\bigl(z_i-z(x)\bigr),
\end{equation}
gdzie $\gamma$ jest estymowany na danych walidacyjnych, a nie przyjmowany
bezkrytycznie jako stały gradient atmosferyczny.

\section{Niepewność}

Dla każdej zmiennej przewidujemy kwantyle, np.
\[
  q_{0.10}(x,u),\quad q_{0.50}(x,u),\quad q_{0.90}(x,u).
\]
Przedział 80\%:
\[
  I_{0.8}(x,u)=
  [q_{0.10}(x,u),q_{0.90}(x,u)].
\]

Warstwa przestrzenna dodaje niepewność wynikającą z odległości od stacji,
liczby sąsiadów i walidacji leave-one-station-out. W uproszczonym wariancie:
\begin{equation}
  \sigma_{\mathrm{total}}^2(x,u)
  =
  \sigma_{\mathrm{model}}^2(x,u)
  +
  \sigma_{\mathrm{spatial}}^2(x,u).
\end{equation}

\section{Aktywacja modelu}

Model nie jest aktywowany tylko dlatego, że został wytrenowany. Dla każdej
zmiennej porównujemy go z modelem bazowym. Przykładowy warunek:
\begin{equation}
  \frac{\operatorname{MAE}_{\mathrm{base}}-
        \operatorname{MAE}_{\mathrm{candidate}}}
       {\operatorname{MAE}_{\mathrm{base}}}
  \ge \eta,
\end{equation}
gdzie $\eta$ jest minimalną wymaganą poprawą.

Dodatkowe warunki:
\begin{itemize}[nosep]
  \item brak regresji jakości dla krytycznych horyzontów;
  \item poprawne pokrycie przedziałów kwantylowych;
  \item brak niefizycznych wartości;
  \item zgodność schematu cech i wersji danych;
  \item przejście testu predykcji kontrolnej.
\end{itemize}

\section{Artefakty i DigitalOcean Spaces}

Wyniki są indeksowane przez dokładny czas:
\begin{lstlisting}[style=code]
hourly/
  models/
    temperature/active.json
    precipitation/active.json
    PM10/active.json
    PM2.5/active.json
  forecasts/
    origin=2026-08-03T10-00-00Z/
      target=2026-08-03T11-00-00Z/
      ...
      target=2026-08-05T10-00-00Z/
  maps/
    origin=.../
      target=.../
        PM10/surface.json.gz
        PM2.5/surface.json.gz
        temperature_c/surface.json.gz
        precipitation_probability/surface.json.gz
        precipitation_mm/surface.json.gz
\end{lstlisting}

App Platform nie wykonuje modelu. Odczytuje dokładny pakiet dla
\texttt{target\_time}. Jeżeli pakiet nie istnieje, zwraca informację o braku
prognozy; nie wybiera po cichu najbliższego czasu.

\section{Informacja prezentowana użytkownikowi}

Dla każdego wyniku interfejs pokazuje:
\begin{itemize}[nosep]
  \item czas bazowy;
  \item dokładny czas docelowy;
  \item horyzont w godzinach;
  \item metodę prognozy;
  \item wersję każdego modelu;
  \item PM10, PM2.5, temperaturę i opad;
  \item prawdopodobieństwo opadu;
  \item przedziały niepewności;
  \item współrzędne punktu zapytania;
  \item współrzędne środka komórki siatki;
  \item odległość pomiędzy punktami;
  \item algorytm interpolacji przestrzennej i liczbę stacji.
\end{itemize}

\section{Dokumentacja dostępna na platformie}

Platforma udostępnia:
\begin{itemize}[nosep]
  \item zakładkę \emph{Model matematyczny};
  \item zakładkę \emph{Jak działa przetwarzanie};
  \item endpoint \texttt{GET /api/v1/docs/mathematics};
  \item endpoint \texttt{GET /api/v1/docs/processing};
  \item pobranie źródła \LaTeX;
  \item opcjonalne pobranie skompilowanego PDF.
\end{itemize}

\section{Wymagania danych}

Bieżący endpoint SYNOP daje najnowszą obserwację, co nie wystarcza do
trenowania modeli pogodowych. Pipeline powinien automatycznie uzupełnić
historię z oficjalnych, miesięcznych archiwów terminowych IMGW. Minimalny
praktyczny zakres pilota:
\begin{itemize}[nosep]
  \item co najmniej 3--6 miesięcy danych godzinowych dla testu technicznego;
  \item pełny rok dla uchwycenia sezonowości;
  \item preferowane 2--3 lata dla stabilniejszej walidacji.
\end{itemize}

Jeżeli historia jest niewystarczająca, system nie może udawać modelu
produkcyjnego. Powinien aktywować jawny fallback i wyświetlać status
\texttt{insufficient\_history}.

\section{Konfiguracja referencyjna}

\begin{lstlisting}[style=code]
hourly_forecasting:
  enabled: true
  min_horizon_hours: 1
  max_horizon_hours: 48
  step_hours: 1
  exact_target_time_required: true
  temporal_interpolation: pchip
  allow_extrapolation: false
  targets:
    - PM10
    - PM2.5
    - temperature_c
    - precipitation_mm
  precipitation:
    occurrence_threshold_mm: 0.1
    model: hurdle
  uncertainty:
    quantiles: [0.10, 0.50, 0.90]
  validation:
    strategy: expanding_window
    cross_fit_weather_for_pm: true
    minimum_folds: 3
\end{lstlisting}

\section{Kryteria odbioru}

Zmiana jest ukończona, gdy:
\begin{enumerate}
  \item prognozy istnieją dla każdej godziny $1,\ldots,48$;
  \item temperatura i opad są pełnoprawnymi celami;
  \item zapytanie o godzinę nie wybiera najbliższego modelu 6/12/24;
  \item modele PM korzystają z prognoz pogodowych uzyskanych bez wycieku;
  \item PCHIP jest używany wyłącznie pomiędzy sąsiednimi godzinami;
  \item opad ma osobny model wystąpienia i wielkości;
  \item metryki są liczone dla każdej zmiennej i każdego horyzontu;
  \item powierzchnie mapy są wersjonowane przez \texttt{target\_time};
  \item API i UI pokazują czas bazowy, docelowy oraz metodę;
  \item dokument matematyczny jest dostępny w aplikacji i jako plik \LaTeX.
\end{enumerate}

\begin{thebibliography}{9}

\bibitem{sklearn-hgbr}
Scikit-learn developers,
\emph{HistGradientBoostingRegressor --- official documentation}.
Dokumentacja opisuje m.in. straty Poissona, gamma oraz kwantylową.

\bibitem{sklearn-tweedie}
Scikit-learn developers,
\emph{TweedieRegressor --- official documentation}.
Uogólniony model liniowy z rozkładem Tweediego.

\bibitem{sklearn-ts}
Scikit-learn developers,
\emph{TimeSeriesSplit --- official documentation}.
Chronologiczna walidacja bez trenowania na przyszłości.

\bibitem{scipy-pchip}
SciPy developers,
\emph{PchipInterpolator --- official documentation}.
Przedziałowa interpolacja Hermite'a zachowująca lokalny kształt.

\bibitem{imgw-open}
IMGW-PIB,
\emph{Dane pomiarowo-obserwacyjne: terminowe dane synoptyczne}.
Oficjalne miesięczne archiwa danych meteorologicznych.

\end{thebibliography}



\section{Ograniczona próba i uczenie przyrostowe}

Niech $\mathcal D$ oznacza kompletne archiwum. Polityka selekcji tworzy
ograniczoną próbę
\begin{equation}
 \mathcal D_{\mathrm{train}}
 = S_{\mathrm{policy}}(\mathcal D;\,p,B),
\end{equation}
gdzie $p$ oznacza profil, a $B$ budżet liczby rekordów i czasu. Selekcja nie
usuwa danych źródłowych.

Waga przykładu:
\begin{equation}
 w_i = 2^{-a_i/\tau}
       n_{s(i)}^{-1/2}
       n_{h(i)}^{-1/2}.
\end{equation}
Zanik eksponencjalny wzmacnia dane najnowsze, natomiast czynniki liczebności
równoważą stacje i horyzonty.

Dla kosztownego modelu bazowego $f_\theta$ definiujemy resztę
\begin{equation}
 r_i=y_i-f_\theta(x_i,h_i).
\end{equation}
Korektor $g_\varphi$ jest aktualizowany przyrostowo, a prognoza ma postać
\begin{equation}
 \widehat y_i=
 \operatorname{clip}\!\left(
 f_\theta(x_i,h_i)+g_\varphi(z_i)
 \right).
\end{equation}
Promocja następuje tylko, jeżeli
\begin{equation}
 \frac{\operatorname{MAE}_{\mathrm{base}}-
       \operatorname{MAE}_{\mathrm{corrected}}}
      {\operatorname{MAE}_{\mathrm{base}}}
 \ge \eta.
\end{equation}

Dla dwóch kolejnych okien zweryfikowanych prognoz:
\begin{equation}
 \Delta_{\mathrm{MAE}}=
 \frac{\operatorname{MAE}_{\mathrm{recent}}-
       \operatorname{MAE}_{\mathrm{reference}}}
      {\operatorname{MAE}_{\mathrm{reference}}}.
\end{equation}
Pełny trening jest rekomendowany, gdy
$\Delta_{\mathrm{MAE}}>\delta_{\mathrm{MAE}}$ albo
$|\operatorname{bias}_{\mathrm{recent}}|>\delta_{\mathrm{bias}}$.


\section{Generyczna rodzina celów jakości powietrza}

Niech $\mathcal P$ będzie konfigurowalnym zbiorem parametrów. Dla
$p\in\mathcal P$ definiujemy
\[
 Y^{(p)}_{s,t,h}=X^{(p)}_{s,t+h}.
\]
Wspólny kontrakt modelu ma postać
\[
 \widehat Y^{(p)}_{s,t,h}=f_{p,\theta_p}
 \left(\bm x^{(p)}_{s,t},h,\bm\phi(t+h),
       \bm w_{s,t+h},\bm z^{(p)}_{s,t}\right),
\]
przy czym każdy parametr może używać innego providera i funkcji straty.
Cechy pomocnicze spełniają warunek kauzalności
$\sigma(\bm z_{s,t})\subseteq\mathcal F_t$.

Jeżeli definicja parametru wyznacza dziedzinę $D_p$, publikowana predykcja jest
projekcją
\[
 \widetilde Y^{(p)}=\Pi_{D_p}(\widehat Y^{(p)}).
\]
Dodanie kodu do katalogu nie powoduje treningu; koszt pojawia się dopiero po
włączeniu roli celu. Selekcja celów pojedynczego przebiegu pozwala trenować
nowy parametr bez dezaktywowania istniejących championów.

\end{document}
